\section{STM32 Firmware Generation Workflow}
\label{sec:firmware}

Workflow 1 automates the initialization of STM32CubeMX projects through CLI-driven scripting, eliminating manual GUI interaction and ensuring reproducible firmware project structure.

\subsection{Project Information Collection}

The \texttt{collect\_project\_info()} node extracts hardware specifications from natural language using structured LLM extraction:

\begin{lstlisting}[language=Python, caption={Firmware project info extraction}]
# User: "Create a project for STM32F401 with STM32CubeIDE, name MyApp"

llm_extractor = ChatTriton(model="mistral").with_structured_output(ProjectInfoExtraction)
result = llm_extractor.invoke([
    SystemMessage(content=project_info_extraction_instructions),
    HumanMessage(content=state.message)
])

# Parsed: {ioc_file_path: null, board_name: "STM32F401", mcu_series: "F4", 
#          project_name: "MyApp", toolchain: "STM32CubeIDE"}
\end{lstlisting}

Alternatively, users can specify an existing \texttt{.ioc} configuration file path, bypassing project creation.

\subsection{STM32Cube Package Installation}

The \texttt{search\_and\_install\_stm32\_package()} function automates the download of the required STM32Cube firmware package from GitHub. This step is critical for headless operations: unlike the STM32CubeIDE/MX GUI, which interactively prompts the user to download missing packages, the CLI mode fails abruptly with a "Project loading error" if the required MCU package is missing from the repository. To resolve this:

\begin{enumerate}
  \item The agent constructs the repository URL based on the MCU series (e.g., \texttt{STM32CubeF4} for F4 series).
  \item It clones the repository directly to \texttt{\textasciitilde/STM32Cube/Repository/}.
  \item It verifies the presence of required HAL files and device headers before proceeding.
\end{enumerate}

This proactive provisioning eliminates the manual "STM32CubeMX Updater" dependency and prevents blocking errors in the automated pipeline.

STM32CubeMX supports headless operation via the \texttt{-s} (script mode) command-line flag. The \texttt{generate\_cubemx\_script()} node creates a dedicated \texttt{.script} file that orchestrates the entire generation process. This script first loads the initial \texttt{.ioc} file (if provided by the user) or falls back to a default board template. It then systematically configures the necessary micro-controller peripherals, such as GPIO, UART, or I2C, based on the specific requirements extracted during the project info collection phase. 

Finally, the script sets the target code generation parameters, including the selected toolchain and the destination output directory, before triggering the final \texttt{generate code} command.

The \texttt{execute\_generation()} node then invokes CubeMX:

\begin{verbatim}
<cubemx_path> -q <script_file.script>
\end{verbatim}

\textbf{Note on Headless Licensing:} A significant deployment challenge verified in the testing environment (using X11 forwarding) is that STM32CubeMX CLI requires a one-time GUI acceptance of the license agreement on first run. Failure to perform this manual initialization (or using \texttt{xvfb-run} without prior acceptance) causes the generation to hang indefinitely. The orchestrated workflow assumes this pre-configuration has been satisfied on the server agent.

\subsection{Overcoming GUI Dependencies: The Pre-seeded IOC Strategy}
\label{subsec:preseeded_ioc}

To resolve the interactive pop-up issues described in Section \ref{subsec:gui_locked_features}, the firmware workflow was refactored to move away from generic board-based loading (\texttt{load <MCU>}) towards a \textbf{Pre-seeded IOC Strategy}.

The agent maintains a repository of pre-validated \texttt{.ioc} templates, one for each supported chip family (e.g., \texttt{STM32H7A3ZITx.ioc}, \texttt{STM32U585AIIxQ.ioc}). These files are configured manually once in the GUI to satisfy the "recommendation wizards" (e.g., MPU regions defined, TrustZone context set to Non-Secure).

The automation logic was subsequently refined to rely on template-based injection. Instead of initializing a raw, unconfigured project, the orchestration agent first identifies the pre-validated template corresponding to the user's specific development board. The generated CubeMX script utilizes a lean scripting approach by executing a direct \texttt{config load "<template.ioc>"} command; critically, by deliberately omitting the standard \texttt{load <MCU>} directive, the tool gracefully skips the architectural negotiation phase and bypasses the interactive, GUI-blocking dialogs. To ensure isolation, project decoupling connects the template to the user's workflow: the \texttt{project name} and \texttt{project path} commands are executed to dynamically re-target the loaded template to the user's specific directory. This guarantees that the final output is a clean, uniquely named project while smoothly inheriting the pre-validated "silent" configuration.

Generated output includes initialization code (\texttt{main.c}, \texttt{stm32f4xx\_hal\_msp.c}), linker scripts, and IDE project files.

\subsection{Human-in-the-Loop Gate}

After successful generation, the \texttt{decide\_continue\_to\_ai()} node presents an interrupt:

\begin{quote}
  \textit{"Firmware project created. Do you want to proceed with AI model selection?"}
\end{quote}

The user's response is classified via LLM, allowing transitions between workflows (firmware $\rightarrow$ AI analysis) or workflow termination.
